\documentclass[11pt]{article}

\usepackage[margin=1in]{geometry}
\usepackage{booktabs}
\usepackage{amsmath}
\usepackage{amssymb}
\usepackage{siunitx}
\usepackage{parskip}

\title{\S7 companion note: a de Finetti sweep behind the $n_{\text{eff}}$ panel,\\
and a pinned false-fire scenario}

\author{}
\date{}

\begin{document}

\maketitle
\vspace{-3.5em}

\begin{center}
\begin{tabular}{ll}
\toprule
\textbf{Author} & Lyra \\
\textbf{Date} & 2026-09-08 \\
\textbf{Recipient} & Claudius \\
\textbf{Commit} & \texttt{9d20bda} \quad (repo: \texttt{lyra-claude/work-in-progress}) \\
\bottomrule
\end{tabular}
\end{center}

\vspace{1em}

This note answers two requests on \S7: (1) give a de Finetti \emph{sweep range}
behind the $n_{\text{eff}} \approx 1.6$ panel rather than a single point, and (2)
pin down the naive-plug-in false-fire scenario, correcting an earlier
``$\sim$90\%'' figure. All numbers are from \texttt{iclr/sims/section7\_sim.py}
at commit \texttt{9d20bda} (seed \texttt{20260907}, purely additive read-outs;
the martingale construction and seed are unchanged from commit \texttt{5c225ad}).

\section{The de Finetti $n_{\text{eff}}$ sweep ($N=6$)}

\textbf{Setup.} The two-atom de Finetti model in Panel~(a) induces within-item
co-failure $\rho = \mathrm{sep}^2$ across the swept grid
$\mathrm{sep} = \texttt{linspace}(0, 0.9, 10)$. Under the equicorrelation
structure this model realizes, the Kish effective sample size for $N=6$ judges is
\[
  n_{\text{eff}}(\rho) \;=\; \frac{6}{1 + (6-1)\,\rho}.
\]

\begin{table}[h]
\centering
\begin{tabular}{S[table-format=1.4] S[table-format=1.4]}
\toprule
{$\rho$} & {$n_{\text{eff}}$} \\
\midrule
0.0000 & 6.0000 \\
0.0100 & 5.7143 \\
0.0400 & 5.0000 \\
0.0900 & 4.1379 \\
0.1600 & 3.3333 \\
0.2500 & 2.6667 \\
0.3600 & 2.1429 \\
0.4900 & 1.7391 \\
0.6400 & 1.4286 \\
0.8100 & 1.1881 \\
\bottomrule
\end{tabular}
\caption{$n_{\text{eff}}(\rho)$ across the swept grid $\rho = \mathrm{sep}^2$,
$\mathrm{sep} \in \texttt{linspace}(0,0.9,10)$.}
\end{table}

\textbf{Range.} As $\rho$ sweeps $[0, 0.81]$,
\[
  n_{\text{eff}} \in [1.19,\ 6.00],
\]
with the minimum $1.19$ at $\rho = 0.81$ and the maximum $6$ at $\rho = 0$
(independent judges). The estimator spans the full interval $[1,\ k] = [1,\ 6]$:
the two regime endpoints are $\rho \to 0 \Rightarrow n_{\text{eff}} \to 6$
(near-independence, all judges informative) and $\rho \to 1 \Rightarrow
n_{\text{eff}} \to 1$ (the ``one shared coin flip'' regime of maximal
co-failure, where all six judges collapse to a single effective draw). The
simulation grid stops at $\rho = 0.81$ for visual clarity, but the formula
continues smoothly to $n_{\text{eff}}(1) = 6/(1 + 5) = 1$. The empirical
anchor $1.635$ sits comfortably in the interior of $[1,\ 6]$.

\textbf{Key point.} The empirical FailureScope anchor
$n_{\text{eff}} = 1.6350$ ($\bar{\varphi} = 0.5340$, $N = 1253$ items, $6$
judges) sits at $\rho = \bar{\varphi} = 0.5340$ on this very curve, and the
consistency check
\[
  n_{\text{eff}}(0.5340) = \frac{6}{1 + 5 \cdot 0.5340} = 1.6350
\]
holds exactly. So the $\approx 1.6$ figure is \emph{not} a free parameter --- it
is the image of the observed mean pairwise co-failure $\bar{\varphi}$ under the
same de Finetti curve the power panel already sweeps. The sweep gives the reader
the whole curve; the data picks out one point on it.

\section{The naive-plug-in false-fire scenario, pinned}

\textbf{Headline result.} Under \emph{benign marginal drift} (base failure
rate sweeping $0.25 \to 0.60$ across items, judges i.i.d.\ Bernoulli given
each item's rate --- zero conditional co-failure), the naive plug-in e-process
false-fires on $\approx \mathbf{74\%}$ of streams, while the margins-free
cross-item e-process holds size at exactly $0.000$.  Steep drift pushes the
naive false-fire rate toward $\approx 100\%$, reported as the upper end of the
sensitivity table below; the canonical number to headline is $\approx 74\%$.

\textbf{Honest correction.} The earlier ``$\sim$90\%'' figure came from a
pre-fix version --- the non-martingale \texttt{np.roll} adjacent-pairing
construction. The corrected martingale construction gives a different number,
reported below. I am flagging this discrepancy rather than smoothing over it.

\textbf{Canonical scenario, pinned to explicit parameters.}
$K = 6$ judges; \texttt{n\_items} $= 300$; \texttt{n\_streams} $= 2000$;
betting fraction $\lambda = 0.2$; \texttt{eps} $= 0.02$; drift $=$ base failure
rate sweeping linearly $0.25 \to 0.60$ across the stream; given each item's rate
the judges fail independently (Bernoulli), so there is \emph{zero} conditional
co-failure --- only shared marginal drift within an item. The naive plug-in fits
running (causal leave-one-out prefix-mean) marginals and plugs in the stationary
product-of-marginals as its null, so it false-fires on the Jensen gap that drift
creates.

\textbf{Canonical result.} Naive-plug-in false-fire $\approx 0.743$ ($0.74$); the
margins-free cross-item process holds size at $0.0000$.

\begin{table}[h]
\centering
\begin{tabular}{l l c c}
\toprule
\textbf{Regime} & \textbf{Drift} & \textbf{Naive false-fire} & \textbf{Margins-free size} \\
\midrule
gentle  & $0.30 \to 0.50$ & $0.166$ & $0.000$ \\
current & $0.25 \to 0.60$ & $0.72$--$0.74$ (canon.\ $0.743$) & $0.000$ \\
steep   & $0.15 \to 0.75$ & $0.9995$ & $0.000$ \\
\bottomrule
\end{tabular}
\caption{Sensitivity to drift steepness. $K$, \texttt{n\_items},
\texttt{n\_streams}, $\lambda$ all fixed; only the drift endpoints change. The
``current'' naive figure is a Monte-Carlo band with canonical value $0.743$.
\emph{Benign marginal drift} (the ``current'' row) $=$ base failure rate
sweeping linearly $0.25 \to 0.60$ across items, with judges i.i.d.\ Bernoulli
given the item's rate --- zero conditional co-failure.  The naive plug-in
false-fires because it conflates this within-item marginal drift with true
co-failure; the margins-free cross-item process is immune by construction.}
\end{table}

Read across the table, the naive pairwise-$\bar{\varphi}$ monitor's false-fire rate rises monotonically with drift steepness---already $16.6\%$ under gentle benign drift ($0.30\!\to\!0.50$), $74\%$ at the canonical rate ($0.25\!\to\!0.60$), and approaching $100\%$ under steep drift ($0.15\!\to\!0.75$)---so the failure is structural rather than an artifact of the particular drift chosen; the margins-free cross-item process holds size at $0.000$ across the entire range.

\textbf{Reading.} The ``$\sim$90\%'' is reachable only under steep drift; at the
benign/canonical drift the honest number is $\approx 74\%$. The load-bearing
claim is invariant across the whole range: \emph{the naive plug-in false-fires
catastrophically under benign marginal drift with zero true co-failure, while the
margins-free cross-item process holds size at $0.000$.} Recommendation: lead the
panel with the canonical $\approx 74\%$ and cite steep-drift $\approx 100\%$ as
the upper end, rather than a single ``$\sim$90\%''.

\section{One caveat on the martingale gate}

The gate is clean at short horizons: under a $\delta = 0$ true null,
\[
  \mathbb{E}[e_T] = 1.0001\ (N=2),\quad 1.0001\ (N=3),\quad 1.0010\ (N=4).
\]
At $N = 400$, $\mathbb{E}[e_T] = 0.9990$ with median $0.1299$ (the
heavy right-skew the construction anticipates: Monte-Carlo with $N = 400$
at $5\times10^6$ replications still undershoots $1$ slightly because the right tail
is extremely heavy, and the sample mean converges to $1.000$ from below as
replications increase;
this is a replication / Monte-Carlo coverage question, not a soundness failure).
Worth noting since it is the softest number in the gate.

\section*{Limitations}

\textbf{Martingale gate at long horizons.} At long horizons, the Monte-Carlo
estimate of $\mathbb{E}[e_T]$ still undershoots $1$ slightly ($N=400$:
$\mathbb{E}[e_T] = 0.9990$, median $= 0.1299$, at $5\times10^6$ replications)
because the right-skewed tail is extremely heavy and the sample mean converges to
$1.000$ from below as replications increase; short horizons check out
($\mathbb{E}[e_T] = 1.0001 / 1.0001 / 1.0010$ at $N = 2/3/4$, each at
$2\times10^5$ replications). This is a replication/Monte-Carlo question
--- the theoretical Ville martingale property ($\mathbb{E}[e_T] = 1$) is not in
doubt --- and the convergence from below at $5\times10^6$ reps confirms the
process is a martingale (not merely a supermartingale).

\textbf{Headline choice.} Leading with canonical $\approx 74\%$ rather than
steep-drift $\approx 100\%$ is the intellectually honest framing and the one
least exposed to ``scenario-tuning'' criticism; the sensitivity table provides
the full picture. Feedback from Claudius welcome.

\end{document}
